% =====================================================================
%  Mapping and Localisation
%  SCOPE: mapping.py (log-odds grid, Bresenham, inflation),
%         localization.py (odometry, likelihood-field scan matching),
%         exploration.py (frontier extraction).
% =====================================================================
\subsection{Mapping and Localisation}
\label{sec:phaseA_mapping}

\subsubsection{The representation: a log-odds occupancy grid}

The map is a $100 \times 100$ grid covering the \SI{10}{m} arena at
\SI{0.10}{m} resolution, following Moravec and Elfes
\cite{moravec1985occupancy}. Each cell carries not a binary state but a
\emph{log-odds} value $\ell$, the logarithm of the ratio between the
probability that the cell is occupied and the probability that it is
free:

\begin{equation}
\ell(c) = \log \frac{P(c\ \text{occupied})}{P(c\ \text{free})}.
\label{eq:logodds}
\end{equation}

The reason for the logarithm is arithmetic, not aesthetic. In
probability space, fusing independent observations requires a product
and a renormalisation; in log-odds space it is an addition
\cite{thrun2005probabilistic}. Fusing a scan therefore becomes a sweep
of \texttt{+=} over an array, which is both fast and numerically
untroubled. Two constants define the sensor model actually used:

\begin{center}\footnotesize
\begin{tabular}{@{}llr@{}}
\toprule
\param{L\_OCC}  & added where a beam terminates & $+0.85$\\
\param{L\_FREE} & added to every cell the beam crosses & $-0.40$\\
\param{L\_CLAMP} & saturation bound on $|\ell|$ & $5.0$\\
\param{L\_OCC\_THRESHOLD} & $\ell$ above which a cell is ``occupied'' & $0.5$\\
\bottomrule
\end{tabular}
\end{center}

Three deliberate asymmetries are encoded in those four numbers.

\textbf{Occupied evidence is worth more than free evidence}
($0.85$ against $0.40$). A beam that stops has made a positive
measurement of a surface; a beam that passes through has only failed to
find one. Making them equal produces a map that erodes thin obstacles,
because a gutter edge is crossed by many more passing beams than
terminating ones.

\textbf{The belief is clamped.} Without \param{L\_CLAMP}, a wall seen
five hundred times reaches $\ell = 400$ and needs a thousand
contradicting observations to be forgotten. The map would then be unable
to react to a moving worker who parks somewhere and later leaves. The
clamp at $\pm 5$ keeps roughly ten contradicting scans sufficient to
flip a cell, which is the right time constant for obstacles moving at
\SI{0.3}{m/s} in a \SI{10}{m} arena.

\textbf{The occupancy threshold is not zero but $+0.5$.} A cell must
accumulate net positive evidence, not merely non-negative evidence,
before the planner is told to avoid it. This is the cheapest available
defence against single-beam noise.

\subsubsection{Integrating one scan: exact ray casting}

For each of the 360 beams the mapper computes the world-frame bearing
$\alpha = \psi + \big(\tfrac{\mathrm{fov}}{2} - i\,\Delta\big)$, decides
whether the beam terminated on a surface ($\texttt{dist}$ finite and
below \param{max\_range}), and walks the integer line from the robot
cell to the endpoint cell with Bresenham's algorithm
\cite{bresenham1965algorithm}:

\begin{lstlisting}[style=python]
for (cc, cr) in self._bresenham(r0c, r0r, ec, er):
    if cc == ec and cr == er:
        break            # stop BEFORE the endpoint
    self.log_odds[cr, cc] += L_FREE
if hit:
    self.log_odds[er, ec] += L_OCC
\end{lstlisting}

The \texttt{break} is not decoration. Without it the endpoint cell
receives \param{L\_FREE} on the same pass that gives it
\param{L\_OCC}, and the wall's net evidence per scan falls from $+0.85$
to $+0.45$ --- halving the convergence rate for no reason at all. It is
a one-line detail that is invisible in the code review and obvious in
the resulting map.

\paragraph{A rejected optimisation, and why it was rejected}
The per-beam Python loop is the slowest part of \phaseA{}. A numpy
rewrite was implemented and benchmarked: sample each ray at half-cell
intervals and scatter with \texttt{np.add.at}. It was measured to be
\emph{both slower and less accurate} --- slower because
\texttt{np.add.at} is not vectorised internally over duplicate indices,
and less accurate because half-cell sampling can skip the very cell a
beam grazes, eroding occupied cells along oblique walls. The exact
sweep was kept, and the mapper was instead run once every
\param{INTEGRATE\_EVERY} $= 4$ control ticks, which costs nothing in map
quality because the robot moves less than \SI{4}{cm} in that interval.
Recording a rejected optimisation is as much part of the engineering
record as recording an accepted one.

\paragraph{A near-range gate}
Beams shorter than \param{min\_range} $= \SI{0.25}{m}$ are discarded
entirely. At that distance the return is either the robot's own
structure or noise, and integrating it stamps a permanent obstacle onto
the cell the robot is standing in --- after which the planner cannot
find a path out of it.

\subsubsection{From belief to plan: thresholding and inflation}

The planner cannot consume log-odds. \texttt{update\_grid\_from\_log\_odds}
produces the binary planning grid in two steps: threshold at
\param{L\_OCC\_THRESHOLD}, then dilate every occupied cell by a disc of
radius \param{inflation}.

Inflation is what turns a map of the \emph{world} into a map of the
\emph{robot's configuration space} \cite{lozanoperez1983spatial}. Once
obstacles are grown by the robot's own radius, the planner may treat the
robot as a dimensionless point, and any path it returns through free
cells is collision-free for the real body. In \phaseA{} the radius is
\SI{0.32}{m}, chosen to cover the YouBot half-diagonal of
$\sqrt{0.29^2 + 0.19^2} = \SI{0.347}{m}$ with a small deliberate
shortfall so that the \SI{1.2}{m} aisles remain traversable.

That choice --- a disc sized on the half-diagonal --- is exactly the
approximation that \phaseB{} later had to abandon. A disc of
\SI{0.347}{m} is correct for a robot that may be rotated arbitrarily,
but it is \SI{16}{cm} too conservative sideways, and in the
\SI{0.80}{m} aisles of the Phase~B greenhouse that conservatism is the
difference between a passable lane and an impassable one. The rectangle
model of Sec.~\ref{sec:safety} replaces it.

The dilation itself is done by shifted whole-array \texttt{OR}s rather
than a per-cell neighbourhood scan:

\begin{lstlisting}[style=python]
for dr in range(-radius, radius + 1):
    for dc in range(-radius, radius + 1):
        if dr*dr + dc*dc > radius*radius:
            continue
        out[r_dst0:r_dst1, c_dst0:c_dst1] |= \
            occupied[r_src0:r_src1, c_src0:c_src1]
\end{lstlisting}

which is $O(\text{radius}^2)$ array operations instead of
$O(\text{cells} \times \text{radius}^2)$ scalar ones --- for a
$100\times100$ grid and a 3-cell radius, 45 array \texttt{OR}s in place
of 450\,000 Python iterations.

% =====================================================================
%  Log-odds occupancy mapping from one lidar scan -- Phase A (mapping.py)
%  Two-column float: the beam loop and the grid pipeline sit side by side.
% =====================================================================
\begin{figure*}[!t]
\centering
\begin{tikzpicture}[fcchain]

% ---------------- INITIALISATION (left column) ------------------------
\node[fcterm] (m0) {\texttt{Mapping} constructed};
\node[fcproc, below=of m0] (m1)
     {allocate $100\times100$ arrays\\
      {\tiny \texttt{log\_odds} float32, \texttt{grid} uint8, res.\ 0.10 m}};
\node[fcproc, below=of m1] (m2)
     {zero the belief\\{\tiny $\ell = 0$ everywhere $=$ ``unknown''}};
\node[fcproc, below=of m2] (m3)
     {fix the sensor model\\
      {\tiny \texttt{L\_OCC}$=+0.85$, \texttt{L\_FREE}$=-0.40$, clamp $\pm5$}};

% ---------------- MAIN LOOP (left column) -----------------------------
\node[fcio, below=9mm of m3] (n0) {scan: 360 ranges + pose $(x,y,\psi)$};
\node[fcdec, below=of n0] (n1) {robot cell inside\\the grid?};
\node[fcfail, left=6mm of n1] (n1b) {skip this scan};
\node[fcproc, below=9mm of n1] (n2)
     {\textbf{for each beam $i$}\\
      $\alpha = \psi + (\mathrm{fov}/2 - i\,\Delta)$};
\node[fcdec, below=8mm of n2] (n3)
     {range $<$ \texttt{min\_range}\\(0.25 m)?};
\node[fcfail, left=6mm of n3] (n3b) {discard\\{\tiny self-hit / noise}};
\node[fcdec, below=9mm of n3] (n4) {finite and\\$<$ \texttt{max\_range}?};
\node[fcproc, below=9mm of n4] (n5)
     {endpoint $= $ hit at \texttt{dist}};
\node[fcproc, right=8mm of n5] (n5b)
     {endpoint $= $ \texttt{max\_range}\\{\tiny free ray, no surface}};

% ---------------- MAIN LOOP (right column) ----------------------------
\node[fcsub, below=9mm of n5] (n6)
     {Bresenham walk to the endpoint\\{\tiny exact integer line}};
\node[fcproc, below=of n6] (n7)
     {$\ell \mathrel{-}= 0.40$ on every crossed cell\\
      {\tiny \textbf{break BEFORE the endpoint}}};
\node[fcdec, below=8mm of n7] (n8) {beam hit a surface?};
\node[fcproc, below=9mm of n8] (n9) {$\ell \mathrel{+}= 0.85$ at the endpoint cell};
\node[fcproc, below=of n9] (n10)
     {clip $\ell$ to $\pm 5$\\{\tiny keeps the map able to forget}};
\node[fcproc, below=of n10] (n11)
     {threshold $\ell > 0.5 \Rightarrow$ occupied};
\node[fcsub, below=of n11] (n12)
     {inflate by 0.32 m (disc)\\{\tiny configuration space, shifted array ORs}};
\node[fcio, below=of n12] (n13) {binary planning grid};

% ---------------- TERMINATION -----------------------------------------
\node[fcproc, below=9mm of n13] (p0)
     {ASCII / PNG dump, compare with supervisor truth};
\node[fcterm, below=of p0] (p1) {controller exit};

% ---------------- edges ------------------------------------------------
\draw[fcline] (m0) -- (m1);
\draw[fcline] (m1) -- (m2);
\draw[fcline] (m2) -- (m3);
\draw[fcline] (m3) -- (n0);
\draw[fcline] (n0) -- (n1);
\draw[fcline] (n1) -- node[fclab,above] {no} (n1b);
\draw[fcline] (n1) -- node[fclab,right] {yes} (n2);
\draw[fcline] (n2) -- (n3);
\draw[fcline] (n3) -- node[fclab,above] {yes} (n3b);
\draw[fcline] (n3) -- node[fclab,right] {no} (n4);
\draw[fcline] (n4) -- node[fclab,right] {yes} (n5);
\draw[fcline] (n4.east) -| node[fclab,pos=0.25,above] {no} (n5b.north);
\draw[fcline] (n5) -- (n6);
\draw[fcline] (n5b.south) |- ($(n6.north)+(0,3mm)$);
\draw[fcline] (n6) -- (n7);
\draw[fcline] (n7) -- (n8);
\draw[fcline] (n8) -- node[fclab,right] {yes} (n9);
\draw[fcline] (n8.east) -- ++(14mm,0)
      node[fclab,pos=0.55,above] {no} |- (n10.east);
\draw[fcline] (n9) -- (n10);
\draw[fcback] (n10.west) -- ++(-10mm,0)
      node[fclab,pos=0.5,above] {next beam} |- (n2.west);
\draw[fcline] (n10) -- (n11);
\draw[fcline] (n11) -- (n12);
\draw[fcline] (n12) -- (n13);
\draw[fcback] (n13.east) -- ++(14mm,0)
      node[fclab,pos=0.5,above] {every 4 ticks} |- (n0.east);
\draw[fcline] (n13) -- (p0);
\draw[fcline] (p0) -- (p1);

% ---------------- phase bands ------------------------------------------
\begin{scope}[on background layer]
  \phaseband{bandinit}{INITIALISATION}{phaseInitLine}{(m0)(m1)(m2)(m3)}{ma}
  \phaseband{bandmain}{MAIN LOOP}{phaseMainLine}
    {(n0)(n1)(n1b)(n2)(n3)(n3b)(n4)(n5)(n5b)(n6)(n7)(n8)(n9)(n10)(n11)(n12)(n13)}{mb}
  \phaseband{bandend}{TERMINATION}{phaseEndLine}{(p0)(p1)}{mc}
\end{scope}

\end{tikzpicture}
\caption{Log-odds occupancy mapping from a single lidar scan
(\phaseA{}, \filename{mapping.py}). The initialisation band allocates
the belief array and fixes the sensor model; nothing in the loop depends
on scene content, so the same code maps the \phaseA{} arena and, ported
unchanged in substance, the \phaseB{} greenhouse. Two details in the
main loop carry disproportionate weight: the Bresenham walk stops
\emph{before} the endpoint, so a wall gains the full $+0.85$ per scan
instead of a net $+0.45$; and the clamp at $\pm5$ is what allows a cell
vacated by a moving worker to be forgotten within about ten scans. The
inflation step at the bottom is the boundary between a map of the world
and a map of the robot's configuration space.}
\label{fig:flow_mapping}
\end{figure*}


\subsubsection{Localisation, part 1: dead reckoning}

\filename{localization.py} first provides a GPS-free pose by integrating
the mecanum forward kinematics of Sec.~\ref{sec:phaseA_locomotion} over
the four wheel encoders. It is exact in the absence of slip and its
error is unbounded and monotone in its presence. Every metre driven adds
error that is never removed, which is the entire justification for what
follows.

\subsubsection{Localisation, part 2: likelihood-field scan matching}

The corrector is a \emph{correlative scan matcher}
\cite{olson2009correlative}: given a pose prior from odometry, search a
small window of $(\Delta x, \Delta y, \Delta\psi)$ offsets for the one
that best explains the current scan against a reference map. The score
of a candidate pose is a likelihood field \cite{thrun2005probabilistic}:
project the scan endpoints into the map, look up each endpoint's
distance $d$ to the nearest obstacle surface, and sum a Gaussian,

\begin{equation}
S(x,y,\psi) = \sum_{k \in \text{beams}}
              \exp\!\left(-\frac{d_k^2}{2\sigma^2}\right),
\qquad \sigma = \SI{0.15}{m}.
\label{eq:likelihood}
\end{equation}

Three implementation decisions carry the performance, and each of them
was made after the naive version failed.

\paragraph{The distance field is seeded from surfaces, not from solids}
The first version computed distance-to-nearest-\emph{occupied-cell}.
This saturates: a scan shifted bodily \emph{inside} a large planting
basin still lands every endpoint on an occupied cell, scores a perfect
$d = 0$ everywhere, and the search has no gradient to climb. Real lidar
only ever sees obstacle \emph{boundaries}, so the breadth-first distance
transform is seeded from occupied cells that touch free space:

\begin{lstlisting}[style=python]
occ = grid > 0
free_neighbor[:-1, :] |= ~occ[1:, :]   # and the three other shifts
return occ & free_neighbor
\end{lstlisting}

Endpoints that drift into an obstacle interior are then penalised, and
the score surface acquires a genuine peak at the true pose.

\paragraph{Coarse-then-fine, with a cached beam layout}
The search runs a coarse pass ($\pm\SI{0.30}{m}$ in steps of
\SI{0.10}{m}, $\pm\SI{0.10}{rad}$ in steps of \SI{0.05}{rad}) followed
by a fine pass around the winner ($\pm\SI{0.06}{m}$ / \SI{0.02}{m},
$\pm\SI{0.04}{rad}$ / \SI{0.02}{rad}). That is $7^3 = 343$ candidates
plus $7^3$ again, rather than the $31^3 = 29\,791$ a single fine sweep
would need. Every beam is subsampled by \param{beam\_stride} $= 6$
(60 of 360 beams), and the subsampled indices and their local bearings
depend only on the scan geometry --- constant across all 686 candidates
--- so they are computed once and cached. The score itself is fully
vectorised, which measured about ten times faster than the equivalent
per-beam loop.

\paragraph{A strictly-better trust gate}
The corrected pose is accepted \emph{only if it scores strictly higher
than the prior}:

\begin{lstlisting}[style=python]
if score <= prior_score:
    return prior, prior_score
\end{lstlisting}

The reason is a failure that was observed rather than predicted. In a
region of long straight walls the score surface has a broad plateau: a
scan shifted along the wall explains the map exactly as well as the true
pose. Without the gate the search settles on an arbitrary point of that
plateau, the result is fed back into the odometry as a reset, and the
next iteration starts from a slightly worse prior --- a slow runaway
driven entirely by ties. Comparing with $\le$ rather than $<$ also
covers the degenerate case of an uninformative scan, where every
candidate scores identically and the search would otherwise adopt a
window-corner offset on no evidence at all.

This single line is, in retrospect, the most important defensive
decision in the localisation stack, and its \phaseB{} descendant is the
rule that a scan-matching estimate may never be trusted merely because
it exists (Sec.~\ref{sec:slam}, where the homemade matcher is shown to
be \emph{worse} than the prior it was meant to improve).

\subsubsection{Assembling the estimator}

\texttt{LocalizationEstimator} bundles the two behind one \texttt{pose()}
callable: odometry integrates every tick, and every
\param{correct\_every} $= 8$ ticks the scan matcher snaps the estimate
back onto the reference map and resets the odometry to the corrected
value. Between corrections the pose tracks the wheels; the drift is
bounded by the correction period rather than by the distance driven.
Because the interface is a single callable, the mission layer can be
switched between ground-truth GPS and this estimator by changing one
argument --- which is how the two were compared quantitatively.

\subsubsection{Frontier extraction: deciding where to look next}

\filename{exploration.py} closes the loop between the map and the
mission. Rather than patrolling a fixed waypoint list, the robot drives
to the boundary between known-free and unknown space
\cite{yamauchi1997frontier}. Working directly on the log-odds array,

\begin{align}
\text{free}(c)    &\iff \ell(c) < -0.5, \\
\text{unknown}(c) &\iff |\ell(c)| < 0.1, \\
\text{frontier}(c)&\iff \text{free}(c) \wedge
   \exists\, c' \in \mathcal{N}_4(c) : \text{unknown}(c').
\end{align}

The three-way test is why the log-odds array is kept rather than
discarded after thresholding: a binary grid cannot distinguish
``observed and empty'' from ``never observed'', and frontier exploration
is exactly the algorithm that needs that distinction.

Frontier cells are clustered by 4-connectivity, clusters smaller than
\param{min\_cluster} $= 4$ cells are dropped as noise, and each surviving
cluster contributes one goal. Two refinements matter:

\begin{itemize}[leftmargin=*,itemsep=1pt,topsep=2pt]
  \item The goal is the frontier cell \emph{nearest the robot}, never
        the cluster centroid. For a ring-shaped frontier --- the normal
        shape early in a run --- the centroid falls in the middle of
        already-explored space, and the robot is sent to a place it has
        already been.
  \item Cells closer than \param{min\_dist} $= \SI{0.6}{m}$ are skipped,
        otherwise the robot twitches in place chasing a frontier it is
        already dissolving.
\end{itemize}

The resulting goals are ordered by a greedy nearest-first tour, and an
empty list is the natural termination condition: no frontier means
everything reachable has been seen. In \phaseB{} this adaptive scheme is
replaced by a deterministic three-lap boustrophedon survey
(Sec.~\ref{sec:phaseB_mission}) --- not because frontier exploration failed,
but because a fixed structured aisle layout makes a coverage pattern
both simpler and time-bounded, and the report's time axis
(Sec.~\ref{sec:results}) made bounded completion time a first-class
requirement.
